\chapter{What would genomic computation add?}
\label{ch:what-adds}

\section*{Begin with a comparison, not a name}
A system stores many possible operations, selects a few in response to context,
and combines their effects. Calling the stored collection a genome can be
suggestive. It does not, by itself, explain how the system differs from an
ordinary program library. This book starts with that uncomfortable comparison
because a useful new abstraction should survive it.

Evolutor is a research program, not a demonstrated general intelligence. Its
questions concern the organization, selection and adaptation of computational
parts. We will distinguish a definition from a mechanism, a mechanism from an
implementation, and an implementation from evidence of an advantage. A
biological analogy may suggest a mechanism; it cannot supply the missing
comparison or experiment.

This chapter develops a deliberately small control example. You will reproduce
two dispatch strategies, prove when they agree, and see why counting selected
operations can hide routing work. No neural-network training is required.
The point is to make the experimental burden visible before adding a large
model whose size could obscure it.

\section{Four operations, one request}
Consider a closed library containing these integer functions:
\[
 f_{\mathrm{double}}(x)=2x,\qquad
 f_{\mathrm{square}}(x)=x^2,\qquad
 f_{\mathrm{absolute}}(x)=|x|,\qquad
 f_{\mathrm{identity}}(x)=x.
\]
A request names one operation and supplies an integer. The request
$(\mathrm{double},-3)$ should produce $-6$. It does not require executing the
other three functions. We may add unused identity operations under distinct
names without changing this request's meaning.

Let the library contain $K$ uniquely named operations. We implement two
strategies. \emph{Indexed dispatch} builds a dictionary once and retrieves the
requested operation by name. \emph{Full-scan dispatch} compares the request
with every stored name, remembers the match and then invokes the selected
operation. The full scan deliberately does not stop early, making its
accounting unambiguous. It is a transparent control, not a claim about the
best implementation of a conditional program.

The distinction between what is stored and what is executed already exists in
this ordinary library. If a proposed genomic system exhibits only that
distinction, our example is a counterexample to the claim that the distinction
alone establishes a novel computational principle. It does not show that all
possible genomic mechanisms are equivalent to dictionary dispatch.

\section{Same semantics, different work}
\begin{definition}[Request semantics]
For a library $L$ with unique names and a declared request $r$ in its domain,
the result on input $x$ is $\operatorname{eval}_{L}(r,x)=f_r(x)$, where $f_r$ is
the operation associated with $r$.
\end{definition}

\begin{proposition}[Dispatch equivalence]
Indexed and full-scan dispatch return the same result when they refer to the
same uniquely named, pure operations, receive the same valid request and input,
and complete normally.
\end{proposition}
\begin{proof}
The index associates $r$ with $f_r$. The scan encounters exactly one matching
name and retains that same $f_r$. Both invoke it once on $x$. Purity and equal
input imply equal results. No equality of routing work follows from this
argument.
\end{proof}

Each premise matters. Duplicate names could make a different construction or
scan policy select a different operation. Mutating the library during a request
could break agreement about its contents. An operation that reads changing
external state need not return the same value in two successive runs. Our
closed arithmetic example excludes these complications so the proposition
states a useful, testable contract rather than a universal assertion.

Table~\ref{tab:dispatch} comes from executing both strategies. Increasing the
stored library from four to 64 operations changes the scan's comparison count
but not its firing count. The indexed implementation records one lookup event
at both sizes. All four outputs are $-6$.

\begin{table}[htbp]
\centering
% Generated by scripts/chapter01_artifacts.py; do not edit.
\begin{tabular}{rlrrrr}
\toprule $K$ & Strategy & Gates & Lookups & Firings & Output \\
\midrule
4 & indexed & 0 & 1 & 1 & -6 \\
4 & full scan & 4 & 0 & 1 & -6 \\
64 & indexed & 0 & 1 & 1 & -6 \\
64 & full scan & 64 & 0 & 1 & -6 \\
\bottomrule\end{tabular}

\caption{Logical event counts for the request $(\mathrm{double},-3)$. A firing
is one invocation of the selected arithmetic function. These are not CPU
instruction counts, elapsed times or learned-model measurements.}
\label{tab:dispatch}
\end{table}

\section{What a cost account must include}
Write a resource record as
\[
 C=(S,B,R,X,T),
\]
where $S$ describes stored representation, $B$ construction or update work,
$R$ routing work, $X$ selected-operation execution and $T$ trace production.
This tuple is a bookkeeping choice for this book, not a new complexity theory.
Its entries need units and a measurement boundary before they can be compared.

For our code, the table reports stored entries and selected logical events.
It does not measure bytes. Constructing the library and its index creates
$O(K)$ entries; that work occurs before dispatch. One dictionary lookup in
the table means one source-level lookup operation, not a guarantee of one
processor instruction or constant worst-case latency. Hashing, collisions,
key representation and memory behavior are outside this event count.

Similarly, a firing of $x\mapsto x^2$ is not a fixed amount of arithmetic work
for integers of arbitrary size. The input's bit length matters. Replacing
small integers with very large integers preserves the semantics while changing
the execution workload. A benchmark must record the input domain as well as
the number of calls.

This explains why ``only one gene fired'' would be an incomplete efficiency
argument. The system might have scored a million gates, loaded a large state,
performed expensive verification or written a detailed trace before and after
that firing. Conversely, a larger persistent library might be worthwhile if
it improves reuse over many requests. We need to measure the tradeoff rather
than deciding it from either the stored size or the active count alone.

\section{The nearest established comparison}
Sparse conditional execution has substantial prior art. The sparsely gated
mixture-of-experts work of \citet{shazeer2017} selects a subset of experts using
a learned gate. Its discussion includes routing, load balancing and practical
execution challenges. Our arithmetic library is not a reproduction of that
neural model; it is a simpler semantic control. Neither source establishes
that an Evolutor design is new or better.

The right comparison depends on the claim. For exact named dispatch, a
dictionary is a serious baseline. For learned input-dependent selection,
learned routing is closer. For discovering reusable compositions, program
search and library-learning approaches become relevant. Comparing a proposed
system only with a deliberately inefficient full scan would leave the most
important alternative untested.

Figure~\ref{fig:dispatch-contract} separates a successful request from the
evidence needed to interpret it. The right-hand lane asks about the library
contract, the cost boundary and the comparison. A successful left-hand lane
does not automatically answer any of those questions.

\begin{figure}[p]
\centering
\includegraphics[width=\textwidth]{../build/figures/evo-g06.pdf}
\caption{EVO-G06. Execution and its evidentiary contract. Solid arrows carry
the request, result or evidence; dashed arrows identify checks needed to
interpret a run. No arrow asserts an advantage over an existing method.}
\label{fig:dispatch-contract}
\end{figure}

\section{How an additional mechanism could earn its place}
A more ambitious system might learn when to reuse an operation, construct a
new composition or revise a library after failure. These possibilities are
research directions, not properties of the present example. Each must be
specified well enough that its absence can be tested.

Suppose a proposed mechanism grows a reusable operation library across tasks.
A concrete question would be whether it reduces the cost of solving held-out
compositions without degrading retained tasks. The mechanism must specify what
can be added, what validation an addition passes, how selection changes and
what state is retained. A name such as ``genome growth'' does not answer those
engineering questions.

\begin{hypothesis}[Illustrative, not yet tested]
Under a fixed task distribution and matched training and inference budgets,
a validated library-growth mechanism reduces held-out composition error
relative to a fixed library and a comparable routing baseline, without
increasing retained-task error beyond a declared tolerance.
\end{hypothesis}

This is a template for a falsifiable study, not a result. Before running it,
choose the tasks, splits, budgets, tolerance, repeated-run protocol and primary
metric. Include a control that spends the same resources without library
growth. If the advantage disappears under that control, a claim about the
growth mechanism is not supported. If retained performance declines beyond
the tolerance, the stated hypothesis fails even if new-task accuracy improves.

There are also less dramatic but worthwhile outcomes. A representation might
make failures easier to localize or modules easier to inspect without improving
predictive accuracy. Such benefits require their own operational tests; they
should not be silently translated into claims of intelligence or asymptotic
speedup. A negative result can sharpen the abstraction by showing which
proposed distinction does no useful work.

\section{Traceability is a design obligation}
Our result record includes the chosen name, output, stored count and event
counts. That helps explain this closed request. It is not a tamper-evident
audit trail or a proof of safe execution. A mutable log can report a plausible
story without establishing that the reported operations occurred.

For a future adaptive system, a reproducible record would also need the library
version, input identity, selection configuration and any stochastic state
needed to replay the request. Claims about \gls{provenance} should state which
of these are captured and what replay checks. Claims about authorization need
a separate policy boundary: identifying the selected operation does not grant
it permission to access an external system.

The current example calls only fixed arithmetic functions. It loads no external
tools, generates no code, changes no permissions and has no online adaptation.
Those exclusions make the example small enough to reason about. They must not
be mistaken for demonstrated safety of a later, more capable architecture.
\index{conditional computation}\index{routing cost}\index{semantic equivalence}

\section{Reproduce and challenge the result}
After installing development dependencies, inspect the four records:
\begin{Verbatim}[fontsize=\small]
from minievolutor.dispatch_reference import chapter_example
for size, method, result in chapter_example():
    print(size, method, result)
\end{Verbatim}
The artifact generator uses these records for the printed table and JSON.
Tests compare both strategies across declared operations and signed inputs,
check arithmetic against independent expressions, and include a large integer.
Unknown requests are rejected rather than silently mapped to identity. These
checks establish behavior of this implementation over tested cases. They do
not measure neural-model quality, wall-clock acceleration or generalization.

\section{Exercises and solution notes}
\begin{enumerate}
\item With $K=1000$, how many comparisons and firings does the full scan record
for a successful request? What important costs remain unmeasured?
\item Modify the scan to stop at the first match. Which proposition remains
valid, and which table entries now depend on operation order?
\item Explain why invoking a function once on a million-bit integer cannot be
equated with invoking it once on a small integer for performance purposes.
\item Give an experimental result that would falsify the library-growth
hypothesis even if held-out composition error improved.
\item An author reports one active module out of a million and concludes that
inference cost is independent of library size. Identify a missing measurement
and an alternative mechanism consistent with the reported observation.
\item State one property that our equivalence proof does not establish about
a learned, stateful or externally connected operation library.
\end{enumerate}

\paragraph{Solution notes.}
(1) 1000 comparisons and one firing; construction, bytes, arithmetic cost and
timing remain unmeasured. (2) Equivalence holds for unique names and the same
pure operations; scan comparisons depend on the match position. (3) Operand
representation and arithmetic work depend on bit length, which firing counts
omit. (4) Retained-task degradation beyond the preregistered tolerance violates
the hypothesis. (5) Measure selection cost; scanning every module's gate can
coexist with one final firing. (6) The proof establishes neither learned
selection quality nor safe effects nor stability under library updates.

\section*{The standard for the rest of the book}
We have established a semantic baseline and exposed a misleading cost shortcut.
We have not established a new theory of computation or an advantage for
Evolutor. The next chapter examines biological inspiration with explicit
limits on the analogy. Later mechanisms must return to this chapter's question:
what changed, compared with which baseline, and what evidence shows it matters?
